\documentclass[11pt]{article}
\usepackage[T1]{fontenc}
\usepackage[utf8]{inputenc}
\usepackage[spanish,es-noshorthands]{babel}
\usepackage{lmodern}
\usepackage[margin=2.3cm]{geometry}
\usepackage{amsmath,amssymb}
\usepackage{booktabs}
\usepackage{xcolor}
\usepackage{enumitem}
\usepackage{graphicx}
\usepackage{parskip}
\usepackage[colorlinks=true,linkcolor=black,urlcolor=blue]{hyperref}

\definecolor{acmuerta}{HTML}{9AA0A6}
\definecolor{aclatente}{HTML}{C9A227}
\definecolor{acactiva}{HTML}{2E7D32}
\definecolor{procest}{HTML}{C62828}
\newcommand{\MUERTA}{\textcolor{acmuerta}{\textsf{MUERTA}}}
\newcommand{\LATENTE}{\textcolor{aclatente}{\textsf{LATENTE}}}
\newcommand{\ACTIVA}{\textcolor{acactiva}{\textsf{ACTIVA}}}
\newcommand{\est}[1]{\textcolor{procest}{\textbf{#1}}}

% Inserta una figura si existe; si no, deja una nota (así compila igual).
\newcommand{\figurasi}[3]{% {archivo}{ancho}{caption}
  \begin{figure}[htbp]\centering
  \IfFileExists{#1}{\includegraphics[width=#2\linewidth]{#1}}%
    {\fbox{\parbox{#2\linewidth}{\centering\vspace{2ex}\itshape Falta la figura
     \texttt{#1}; generarla con su script antes de compilar.\vspace{2ex}}}}%
  \caption{#3}\end{figure}}

\title{\textbf{AstroBioSim — Presentación final del proyecto}\\[3pt]
\large Decisiones fundamentales y hallazgos, integrados por las tres áreas}
\author{Equipo AstroBioSim — Matemática $\cdot$ Física $\cdot$ Biotecnología}
\date{Julio 2026}

\begin{document}
\maketitle

\begin{abstract}
\noindent AstroBioSim simula la \textbf{supervivencia poblacional de
microorganismos} en subsuelos de Tierra, Marte y Encelado mediante un
\textbf{autómata celular} estocástico y espacialmente explícito sobre una grilla
2D, gobernado por tres variables ambientales: temperatura $T$, radiación UV $R$ y
actividad de agua $a_w$. Este documento reúne, en un solo hilo, lo que debe
contarse en la exposición final: las \emph{decisiones fundamentales} de cada área
y los \emph{hallazgos} del modelo —el umbral crítico de microrefugios, la solidez
estadística de las corridas y el análisis de sensibilidad que blinda el resultado.
Los tres documentos por disciplina (\texttt{matematicas.tex}, \texttt{fisica.tex},
\texttt{biotecnologia.tex}) desarrollan cada pieza en detalle; este las integra.
\end{abstract}

\section{El proyecto en una frase, y la pregunta de investigación}

Modelamos si una población microbiana \emph{persiste o se extingue} en un subsuelo
planetario, y qué hace falta para que persista. El resultado que ancla todo el
proyecto es que \textbf{el regolito marciano en bulk es inhabitable} para el
crecimiento activo ($a_w\approx 0{,}38$, muy por debajo del límite de división
celular $\approx 0{,}605$): la vida solo puede persistir en \emph{microrefugios}
húmedos transitorios (salmueras delicuescentes). De ahí la pregunta de
investigación (ADR-0015):

\begin{quote}\itshape
¿Con qué frecuencia y magnitud mínimas deben aparecer microrefugios húmedos
transitorios para que una población persista en un subsuelo planetario, en vez de
extinguirse?
\end{quote}

\noindent Que Marte dé \emph{extinción del crecimiento activo} no es un error del
modelo: es el resultado correcto y el punto de partida de la pregunta.

\section{El modelo: tres decisiones fundamentales, una por área}

\paragraph{Matemática — un autómata celular de \emph{tres} estados con cinética
continua.} A diferencia del Juego de la Vida (vivo/muerto), el espacio de estados
es $\{\MUERTA,\LATENTE,\ACTIVA\}$ (ADR-0012): \MUERTA{} es absorbente (la muerte es
irreversible; repoblar es \emph{colonización}, no resurrección), \LATENTE{} está
viva sin reproducirse, \ACTIVA{} crece. La habitabilidad deja de ser binaria: una
\textbf{tasa de crecimiento} continua (gamma concept + CTMI, ADR-0013)
$\mu=\mu_{\mathrm{opt}}\,\gamma_T\,\gamma_{a_w}\,\gamma_{\mathrm{UV}}$ fija la
probabilidad de reproducción por tick ($\Delta t = 1$~h). La actualización es
\textbf{síncrona} (doble buffer, vectorizada en NumPy) y la regla de transición es
\textbf{intercambiable} (ADR-0016) con un editor por bloques y panel de notación
formal (ADR-0018). Un orquestador único corre el bucle bajo dos modos (Sandbox y
Analógico, ADR-0017), y toda la aleatoriedad usa un generador \textbf{sembrado}
(reproducibilidad exacta).

\paragraph{Física — un campo de tres variables que \emph{discriminan}.} El criterio
rector es que una variable que no separa escenarios no entra al modelo. Por eso se
eliminó la presión (ADR-0008, casi constante en los análogos) y la radiación se
reencuadró como \textbf{irradiancia UV}, no dosis ionizante (ADR-0014): la dosis
ionizante marciana ($0{,}077$~Gy/año) tardaría $\sim\!65\,000$ años en alcanzar la
letal de \emph{D.\ radiodurans}, mientras que el UV esteriliza de inmediato. Cada
subsuelo tiene su modelo de campo: Marte amortigua la onda térmica y el UV con la
profundidad (el UV cae mucho más rápido que el calor), Encelado suma fumarolas
gaussianas sobre un fondo frío. Los eventos estocásticos perturban \emph{solo} el
campo, y deben ser \textbf{simétricos}: junto a los que degradan (micro-fisura)
existe el que mejora (salmuera), sin el cual la pregunta de investigación no
tendría mecanismo.

\paragraph{Biotecnología — tres especies análogas con umbrales trazables.} Un
control mesófilo (\emph{E.\ coli}, Tierra), un candidato radiorresistente y
anhidrobiótico (\emph{D.\ radiodurans}, Marte) y una metanógena psicrotolerante
(\emph{M.\ burtonii}, Encelado; ADR-0011). La decisión de fondo es separar
\textbf{crecimiento} ($\mu>0$) de \textbf{supervivencia} ($\mu=0$, sigue viva): es
lo que permite la \emph{latencia} (anhidrobiosis), sin la cual habría que matar a
\emph{D.\ radiodurans} al secarse (falso) o declarar crecimiento bajo el límite de
la vida (indefendible). Todo parámetro lleva \textbf{procedencia}: \textbf{[LIT]}
publicado, \textbf{[DER]} derivado, \textbf{[ANA]} por analogía, \textbf{[CONV]}
convención, \est{[EST]} estimación (deuda). Un parámetro sin procedencia no entra
al modelo.

\section{Resultado base: cada especie en su entorno}

Con los parámetros calibrados y \emph{sin ajustar nada}, cada especie análoga en su
entorno produce:

\begin{center}\small
\begin{tabular}{@{}lcccl@{}}
\toprule
Especie / entorno & \ACTIVA & \LATENTE & \MUERTA & Lectura \\
\midrule
\emph{E.\ coli} / Tierra & $100\%$ & $0\%$ & $0\%$ & el control prospera \\
\emph{D.\ radiodurans} / Marte & $0\%$ & $98\%$ & $2\%$ & \textbf{sobrevive latente, no crece} \\
\emph{M.\ burtonii} / Encelado & $100\%$ & $0\%$ & $0\%$ & crece hasta en los núcleos de fumarola \\
\bottomrule
\end{tabular}
\end{center}

En Marte, combinar la atenuación rápida del UV con la lenta del calor hace emerger
una \textbf{banda de profundidad}: la superficie inmediata queda \MUERTA{} (UV
sobre la fluencia letal) y el resto \LATENTE{} (el UV ya no mata, pero falta agua
para crecer). La ventana habitable queda acotada \emph{arriba} por el UV y
\emph{abajo} por el agua. \textbf{Validación independiente:} el modelo da
$\mathrm{UV}_{\mathrm{sup}}\approx 42{,}2$~W/m$^2$, dentro del rango publicado de UV
marciano ($42$--$55$~W/m$^2$), por una vía distinta a la que lo fijó.

\section{Hallazgo central: el umbral crítico de microrefugios (ADR-0015)}

Operativamente, la pregunta de investigación es un \textbf{barrido}
$\text{frecuencia}\times\text{magnitud}$ del evento de salmuera sobre ensambles
Montecarlo, con la persistencia como variable de respuesta. El hallazgo tiene dos
capas:

\begin{itemize}[nosep]
  \item \textbf{La extremófila no necesita refugios para \emph{sobrevivir}.}
  \emph{D.\ radiodurans} tiene $a_{w,\text{sup}}=0$ (anhidrobiosis total): en Marte
  persiste \LATENTE{} indefinidamente sin ningún refugio, así que con la métrica
  ``viva'' \emph{no hay} umbral crítico —es un hallazgo, no un fallo. El umbral
  interesante para ella es el de \emph{reactivación} (¿los refugios la vuelven a
  \ACTIVA?).
  \item \textbf{La mesófila sí, y con exigencia cuantificable.} \emph{E.\ coli}
  muere sin agua ($a_{w,\text{sup}}=0{,}5$), de modo que da un umbral limpio: para
  persistir en Marte necesita microrefugios \textbf{grandes} (radio $\gtrsim 10$--$16$
  celdas) \textbf{y} \textbf{frecuentes} ($\gtrsim 0{,}27$ disparos/tick). Por debajo
  de ese contorno se extingue; por encima, persiste.
\end{itemize}

\figurasi{barrido_microrefugios_marte_ecoli_radio_viva.png}{0.86}{Barrido
frecuencia $\times$ radio del microrefugio para \emph{E.\ coli} en Marte. El
\emph{umbral crítico} (contorno donde la probabilidad de persistencia cruza $0{,}5$)
separa la extinción de la persistencia: hace falta un refugio a la vez amplio y
frecuente. Reproducible con \texttt{scripts/barrido\_microrefugios.py}.}

\section{Solidez estadística: ¿cuántas réplicas hacen falta? (validación Montecarlo)}

Como el sistema es estocástico, \textbf{una corrida no es la respuesta}. Se
promedian $N$ réplicas con semillas explícitas y se reporta la
\textbf{media $\pm$ desviación muestral}, más el error estándar
$\mathrm{SE}=\sigma/\sqrt{N}$. El estudio de convergencia responde los tres
criterios de validación:

\begin{itemize}[nosep]
  \item \textbf{$N$ suficiente:} cerca del umbral crítico (donde más importa),
  hacen falta $\mathbf{N=180}$ réplicas para que $\mathrm{SE}\le 3\%$ —bastante más
  que las $30$ de una corrida ingenua. La media se estabiliza y el SE decae
  exactamente como $1/\sqrt{N}$.
  \item \textbf{No confundir una corrida con la distribución:} se reporta siempre
  media $\pm\,\sigma$ (p.\ ej.\ $0{,}212\pm 0{,}410$ para la persistencia).
  \item \textbf{Reproducibilidad:} misma semilla $\Rightarrow$ misma traza, exacta.
\end{itemize}

\figurasi{validacion_montecarlo_marte_ecoli_persistencia.png}{0.92}{Convergencia del
estimador Montecarlo. Izquierda: la media acumulada se estabiliza con banda de
confianza del $95\%$. Derecha: el error estándar decae como $1/\sqrt{N}$ y cruza el
objetivo del $3\%$ en $N\approx 180$. Reproducible con
\texttt{scripts/validacion\_montecarlo.py}.}

\section{Blindaje: análisis de sensibilidad de los umbrales (ADR-0013)}

¿Qué tan frágil es la conclusión frente a la incertidumbre de los umbrales
biológicos? Se varía cada umbral sobre su rango de incertidumbre —cuyo ancho sale
de su \emph{procedencia}: [LIT] estrecho, [ANA] el rango documentado
$2{,}5$--$13{,}8\times$ de \emph{M.\ burtonii}, \est{[EST]} ancho— dejando el resto
en su valor nominal, y se mide cuánto se mueve la persistencia. El hallazgo es
contundente:

\begin{quote}
\textbf{Domina un solo umbral: $\boldsymbol{a_{w,\text{sup}}}$ de \emph{E.\ coli}},
que es \est{[EST]} (estimado, sin cita). Barrerlo en su rango mueve la persistencia
$0{,}83$ (de $\sim\!0{,}10$ a $\sim\!0{,}93$) y \textbf{cruza la frontera $0{,}5$};
todos los demás umbrales (temperatura, UV, $a_{w,\min}$) son \emph{planos}. La
conclusión de persistencia \textbf{no es robusta}: pende por completo del parámetro
menos defendible del modelo.
\end{quote}

\noindent Esto no debilita el proyecto: lo \emph{fortalece}. El análisis identifica
exactamente dónde invertir para que el resultado sea defendible —cerrar la deuda
\est{[EST]} de $a_{w,\text{sup}}$— en vez de dejarlo a la intuición.

\figurasi{sensibilidad_marte_ecoli_viva.png}{0.9}{Diagrama de tornado: cada barra es
el rango de persistencia al barrer un umbral en su incertidumbre, ordenadas por
dominancia y coloreadas por procedencia. Una sola barra ($a_{w,\text{sup}}$, \est{[EST]})
cruza la frontera $0{,}5$; el resto no mueve la aguja. Reproducible con
\texttt{scripts/sensibilidad\_umbrales.py}.}

\section{Honestidad científica: deuda abierta}

Parte de la defensa es declarar lo que falta. Estado a la fecha:

\begin{itemize}[nosep]
  \item \textbf{Resuelto:} la $a_w$ media de Marte se re-derivó de la humedad
  relativa cruda ($0{,}382$, antes se usaba el mínimo diario $0{,}19$); el $a_w$ del
  control terrestre se corrigió a $0{,}99$ (era humedad del aire); y
  $\tau=\texttt{SEGUNDOS\_UV\_POR\_TICK}$ se ajustó a $1$~h para coincidir con el
  $\Delta t$ del autómata (antes $8$~h dejaban los umbrales UV $8\times$ permisivos).
  \item \textbf{Abierto — física:} la asíntota térmica de Marte usa la media de los
  mínimos diarios ($7{,}8\,^\circ$C), pero una onda amortigua hacia la \emph{media
  anual} ($\approx 22\,^\circ$C): revisar. Y el pico $\Delta T=25$ de la emisión
  hidrotermal, apilado sobre una fumarola, lleva $T$ a $\sim\!59\,^\circ$C (mata a
  \emph{M.\ burtonii} en el $29\%$ de los disparos): puede ser correcto, pero hay
  que decidirlo.
  \item \textbf{Abierto — biología, y ahora \emph{prioritario}:} los umbrales de
  \emph{supervivencia} $t_{\text{sup}}$ y $a_{w,\text{sup}}$ (de \emph{E.\ coli} y
  \emph{M.\ burtonii}) son \est{[EST]} sin cita directa. El análisis de sensibilidad
  (\S7) muestra que $a_{w,\text{sup}}$ de \emph{E.\ coli} es \textbf{el} parámetro que
  decide el resultado: cerrarlo con literatura es la mejora de mayor impacto pendiente.
\end{itemize}

\section*{Qué contar en la exposición (síntesis)}

\begin{center}\small
\begin{tabular}{@{}llll@{}}
\toprule
\# & Bloque & Mensaje clave & Apoyo \\
\midrule
1 & Pregunta & Marte en bulk es inhabitable; ¿qué refugios salvan a la vida? & --- \\
2 & Modelo & AC de 3 estados + cinética CTMI + campo UV/agua + 3 especies & \texttt{cinetica\_ctmi} \\
3 & Resultado base & control prospera; \emph{D.\ rad.}\ latente en Marte; banda de profundidad & tabla \\
4 & Hallazgo central & umbral crítico: refugio amplio \emph{y} frecuente para \emph{E.\ coli} & fig.\ barrido \\
5 & Rigor & $N=180$ réplicas para $\mathrm{SE}\le 3\%$; media $\pm\,\sigma$; reproducible & fig.\ validación \\
6 & Blindaje & todo pende de $a_{w,\text{sup}}$ \est{[EST]}: dónde mejorar & fig.\ sensibilidad \\
7 & Honestidad & deuda abierta declarada (física + $a_{w,\text{sup}}$) & --- \\
\bottomrule
\end{tabular}
\end{center}

\noindent\textbf{Aporte del proyecto.} La combinación \emph{espacialidad $+$
estocasticidad $+$ latencia}, con umbrales trazables y un análisis de sensibilidad
que dice de qué depende el resultado, es un nicho poco cubierto por la literatura de
habitabilidad —que suele ser puntual (0-D) y determinista. No afirmamos que haya
vida en Marte: cuantificamos \emph{qué condiciones de refugio la harían persistir},
y con cuánta confianza podemos decirlo.

\end{document}
